\chapter{The Interface, Navigation and the Document}
\label{ch:interface}

This chapter is a fast orientation to the FreeCAD window and the handful
of interactions you will repeat thousands of times.

\section{Anatomy of the window}

\begin{description}
  \item[Workbench selector.] The dropdown (or the \emph{View
        $\rightarrow$ Workbench} menu) switches the active toolset. Its
        contents are context-specific to the task, not the document.
  \item[Tree view (Model tab).] The document object graph
        (\cref{ch:introduction}). Objects nest under their consumers;
        the tip feature of a solid body sits at the bottom.
  \item[Property editor.] Below the tree; shows the properties of the
        selected object, split into \emph{Data} (geometric/parametric)
        and \emph{View} (appearance). Editing a Data property and
        recomputing changes the model.
  \item[3D view.] The OpenGL viewport. Selection here is synchronised
        with the tree.
  \item[Report view / Python console.] Optional panels (\emph{View
        $\rightarrow$ Panels}) that surface messages and give a live
        Python prompt.
\end{description}

\section{Navigation}

FreeCAD offers several \emph{navigation styles} (Edit $\rightarrow$
Preferences $\rightarrow$ Display $\rightarrow$ Navigation, or the
selector in the status bar). The default \emph{CAD} style is:

\begin{center}
\begin{tabularx}{\linewidth}{lX}
  \toprule
  Action & Input \\
  \midrule
  Pan    & middle mouse button + drag \\
  Zoom   & scroll wheel \\
  Rotate & middle button + right button (or the on-screen navigation cube) \\
  Select & left click \\
  \bottomrule
\end{tabularx}
\end{center}

The \textbf{navigation cube} in the corner gives named standard views
(Front, Top, Isometric). The numeric keys \texttt{1}--\texttt{6} snap to
the standard orthographic views; \texttt{0} gives the isometric view.

\section{Units and preferences}

Set your unit system early: \emph{Edit $\rightarrow$ Preferences
$\rightarrow$ General $\rightarrow$ Units}. FreeCAD stores lengths
internally in millimetres regardless of the display unit. For FEA this
matters enormously: material properties, forces and results must be in a
\emph{consistent} unit system, a point we return to with emphasis in
\cref{ch:fem-workbench}.

\begin{fcwarn}[Consistent units]
FreeCAD's default internal length unit is the millimetre. If a Young's
modulus is entered in \si{\mega\pascal} ($=\si{\newton\per\milli\metre\squared}$),
lengths in \si{\milli\metre} and forces in \si{\newton}, then stresses
come out in \si{\mega\pascal}. Mixing this with SI-metre inputs is the
single most common source of nonsensical FEA results.
\end{fcwarn}

\section{Saving and the file format}

An \texttt{.FCStd} file is a ZIP archive containing the serialised
object graph, the OCCT geometry (BREP), a thumbnail and GUI state.
Because it is just a ZIP, models are diff-unfriendly but robust and
self-contained. FreeCAD keeps an undo/redo transaction stack per
document; each GUI command is one transaction.
